\uuid{dc1L}
\titre{Zéros de fonction}
\auteur{bodin ; hanani ; mzari}
\organisation{exo7}
\licence{CC BY-NC-SA 4.0}
\datecreate{2020-01-20}
\niveau{L1}
\module{Analyse}
\chapitre{Continuité, limite et étude de fonctions réelles}
\sousChapitre{Continuité : théorie}
\difficulte{2}
\theme{Théorèmes des valeurs intermédiaires, zéros de fonction}
\exo7qcmid{lille-1.continuite.021}

\contenu{
\begin{question}
\qclassification{123.01,123.02}
\qtags{motcle=zéros de fonction}
Soit $f : [a,b] \to \Rr$ une fonction continue (avec $a < b$). Quelles assertions sont une conséquence du théorème des valeurs intermédiaires ?
\begin{answers}
    \bad{Si $f(a) \cdot f(b) > 0$ alors $f$ s'annule sur $[a,b]$.}

    \good{Si $f(a) < k < f(b)$ alors $f(x)-k$ s'annule sur $[a,b]$.}

    \bad{Pour $I \subset \Rr$, si $f(I)$ est un intervalle alors $I$ est un intervalle.}

    \bad{Si $c \in ]a,b[$ alors $f(c) \in ]f(a),f(b)[$.}
\end{answers}
\begin{explanations}
Les trois façons d'énoncer le théorème des valeurs intermédiaires, pour $f : [a,b] \to \Rr$ continue :
(1) si $f(a) \cdot f(b) \le 0$ alors $f$ s'annule sur $[a,b]$ ;
(2) si $f(a) < k < f(b)$ alors il existe $a < c < b$ tel que $f(c)=k$ ;
(3) si $I$ est un intervalle, alors $f(I)$ est un intervalle.
\end{explanations}
\end{question}
}
